%!TEX root = ../main.tex
\chapter{Technologia blockchain}
\label{ch:blockchain}

\section{Rozproszony rejestr i niezmienność}
\label{sec:rozproszony-rejestr}

Blockchain, wprowadzony jako fundament systemu Bitcoin \cite{nakamoto2008bitcoin}, jest w istocie rozproszonym rejestrem (\emph{distributed ledger}): uporządkowanym dziennikiem transakcji replikowanym przez zbiór wzajemnie nieufających sobie węzłów, które za pomocą protokołu konsensusu uzgadniają jedną, wspólną historię. Z perspektywy teorii systemów rozproszonych jest to realizacja replikowanej maszyny stanów (\emph{state machine replication}): jeżeli wszystkie węzły rozpoczynają od identycznego stanu początkowego i wykonują te same transakcje w tej samej kolejności, a logika przetwarzania jest deterministyczna, to stan każdego uczciwego węzła pozostaje identyczny. Uzgodnieniu podlega więc wyłącznie \emph{kolejność} transakcji. Stan jest jej deterministyczną konsekwencją.

Transakcje grupowane są w bloki, a każdy blok zawiera skrót kryptograficzny swojego poprzednika. Pozycję bloku w tym ciągu, liczoną od bloku początkowego (\emph{genesis}), nazywa się jego \textbf{wysokością} (\emph{height}); wysokość łańcucha to numer ostatniego zatwierdzonego bloku i zarazem miara postępu konsensusu, której wzrost oznacza, że sieć finalizuje kolejne bloki. Powstały łańcuch wiąże całą historię: modyfikacja dowolnej zatwierdzonej transakcji zmienia skrót jej bloku, a przez to skróty wszystkich bloków następnych, czyniąc manipulację natychmiast wykrywalną. Ta właściwość, określana mianem niezmienności (\emph{immutability}), ma charakter warunkowy. Łańcuch skrótów gwarantuje wykrywalność modyfikacji, natomiast to protokół konsensusu decyduje, czy alternatywna historia może zostać uznana za obowiązującą (sekcja~\ref{sec:konsensus}).

Istotne dla dalszych rozważań jest rozróżnienie dwóch modeli uczestnictwa. W sieciach otwartych (\emph{permissionless}), takich jak Bitcoin czy Ethereum, każdy może uczestniczyć w konsensusie, a odporność na ataki Sybil zapewnia koszt ekonomiczny (obliczenia lub stawka). W sieciach konsorcyjnych (\emph{permissioned}) zbiór walidatorów jest jawny i kontrolowany, co pozwala stosować klasyczne protokoły konsensusu bizantyjskiego o silniejszych gwarancjach. Rejestr certyfikatów dla Web PKI naturalnie pasuje do modelu konsorcyjnego: kandydatami na walidatorów są podmioty już dziś tworzące infrastrukturę zaufania w Internecie (operatorzy rejestrów domen, dostawcy przeglądarek, organizacje standaryzacyjne, dotychczasowe CA), a wartością dodaną względem stanu obecnego jest nie anonimowość uczestników, lecz brak możliwości jednostronnego działania któregokolwiek z nich.

\section{Funkcje skrótu i drzewa Merkle}
\label{sec:merkle}

Kryptograficzna funkcja skrótu odwzorowuje dane dowolnej długości na skrót stałej długości w sposób praktycznie nieodwracalny (odporność na odtworzenie przeciwobrazu) i bezkolizyjny (niewykonalność znalezienia dwóch wejść o tym samym skrócie). W niniejszej pracy, podobnie jak w większości systemów blockchain, wykorzystywana jest funkcja SHA-256. Skrót pełni rolę zwięzłego, wiążącego zobowiązania (\emph{commitment}) wobec danych: kto zna skrót, może zweryfikować integralność danych dowolnego rozmiaru, porównując jedną 32-bajtową wartość.

Drzewo Merkle \cite{merkle1987digital} rozszerza tę ideę ze zbioru danych traktowanego w całości na jego elementy. Liście drzewa zawierają skróty poszczególnych elementów, a każdy węzeł wewnętrzny zawiera skrót konkatenacji swoich dzieci; skrót korzenia stanowi zobowiązanie wobec całej zawartości. Kluczową właściwością jest możliwość konstruowania zwięzłych dowodów włączenia (\emph{inclusion proof}): aby przekonać weryfikującego, że element należy do zbioru o danym korzeniu, wystarczy przedstawić skróty węzłów-rodzeństwa na ścieżce od liścia do korzenia. Dowód ma rozmiar $O(\log N)$ i jest weryfikowalny bez dostępu do pozostałych danych. To dokładnie ta konstrukcja umożliwia działanie logów Certificate Transparency omówionych w sekcji~\ref{sec:ct}.

Klasyczne drzewo Merkle uwierzytelnia statyczną listę; rejestr certyfikatów wymaga natomiast uwierzytelnionego słownika klucz--wartość, który podlega modyfikacjom i w którym dowieść trzeba także \emph{nieobecności} klucza (pytanie ,,czy domena ma zarejestrowany certyfikat?'' musi mieć weryfikowalną odpowiedź negatywną). Potrzebę tę realizuje rzadkie drzewo Merkle (SMT) \cite{dahlberg2016smt}: drzewo binarne o stałej głębokości równej długości skrótu klucza (256 poziomów dla SHA-256), w którym pozycja liścia jest wyznaczona bitami skrótu klucza, a przygniatająca większość poddrzew jest pusta i reprezentowana przez znane z góry skróty domyślne. Dzięki temu, mimo astronomicznej przestrzeni adresowej $2^{256}$, przechowywane i przeliczane są wyłącznie ścieżki kluczy faktycznie obecnych, a dowód nieobecności ma tę samą postać co dowód włączenia, jest bowiem ścieżką prowadzącą do poddrzewa pustego. Rzadkie drzewo Merkle stanowi rdzeń warstwy stanu projektowanego systemu (rozdział~\ref{ch:projekt}), a jego korzeń, publikowany w każdym bloku, pełni rolę korzenia zaufania, względem którego klienci weryfikują odpowiedzi o stan dowolnej domeny.

\section{Mechanizmy konsensusu}
\label{sec:konsensus}

Protokół konsensusu rozstrzyga, który blok zostaje dołączony do łańcucha, a więc - w języku replikowanej maszyny stanów - ustala wspólną kolejność transakcji pomimo awarii i złośliwego zachowania części uczestników. Wybór protokołu determinuje dwa parametry o pierwszorzędnym znaczeniu dla rejestru certyfikatów: model finalności (czy zatwierdzony blok może zostać wycofany) oraz próg odporności (jaką część uczestników musi kontrolować napastnik, aby naruszyć gwarancje systemu).

\subsection{Klasyczne podejścia: dowód pracy i dowód stawki}
\label{subsec:pow-pos}

W konsensusie opartym na dowodzie pracy (PoW), zastosowanym w Bitcoinie \cite{nakamoto2008bitcoin}, prawo zaproponowania bloku uzyskuje węzeł, który jako pierwszy rozwiąże kosztowną obliczeniowo zagadkę, a obowiązującą historią jest łańcuch o największym skumulowanym nakładzie pracy. Konstrukcja ta dopuszcza chwilowe rozwidlenia (\emph{forks}): dwa bloki wydobyte niemal równocześnie konkurują ze sobą, a jeden z nich zostaje ostatecznie porzucony wraz z zawartymi w nim transakcjami. Finalność ma więc charakter wyłącznie probabilistyczny. Nie istnieje moment, w którym zapis staje się formalnie nieodwracalny; prawdopodobieństwo wycofania bloku maleje wykładniczo z liczbą bloków dobudowanych nad nim, lecz nigdy nie osiąga zera. W praktyce przyjmuje się zatem umowny próg liczby potwierdzeń, po którym zapis uznaje się za dostatecznie pewny, na przykład sześć kolejnych bloków w Bitcoinie, co odpowiada około godzinie oczekiwania.

Dowód stawki (PoS) \cite{king2012ppcoin} zastępuje koszt obliczeniowy kosztem kapitałowym: prawo proponowania bloków przydzielane jest proporcjonalnie do zdeponowanej stawki, której utrata grozi za udowodnione naruszenie protokołu. Eliminuje to charakterystyczne dla PoW zużycie energii, jednak warianty oparte na regule najdłuższego łańcucha oferują jedynie probabilistyczne (a nie ostateczne) potwierdzanie transakcji, a ponadto wprowadzają specyficzne wektory ataku, w tym ataki dalekiego zasięgu (\emph{long-range attacks}), w których napastnik dysponujący historycznymi kluczami walidatorów (np.\ wykupionymi od uczestników, którzy dawno wycofali stawkę) odtwarza alternatywną historię łańcucha od odległego punktu w przeszłości.

Dla rejestru certyfikatów obie te właściwości są dyskwalifikujące w stopniu większym niż dla rejestru płatności. Wycofanie bloku oznaczałoby, że certyfikat uznany za ważny (lub unieważniony) przestaje taki być wstecznie, a klient mógł już podjąć na jego podstawie decyzję o zaufaniu połączeniu. Z kolei atak dalekiego zasięgu pozwalałby fabrykować całe alternatywne historie rejestru na potrzeby ofiar pozbawionych aktualnego widoku sieci. Rejestr certyfikatów wymaga finalności deterministycznej: blok raz zatwierdzony musi pozostać zatwierdzony na zawsze.

\subsection{Konsensus bizantyjski (BFT) i Tendermint}
\label{subsec:bft-tendermint}

Problem osiągania zgody w obecności uczestników zachowujących się dowolnie złośliwie sformalizowali Lamport, Shostak i Pease jako problem generałów bizantyjskich \cite{lamport1982byzantine}, wykazując, że rozwiązanie wymaga, by uczestników złośliwych było mniej niż jedna trzecia: $n \geq 3f + 1$ dla $f$ uszkodzonych węzłów. Granica ta jest fundamentalna i niezależna od konkretnego protokołu. Praktyczną realizację dla systemów asynchronicznych zapoczątkował protokół PBFT \cite{castro1999pbft}, w którym repliki uzgadniają kolejne decyzje w rundach głosowań, a każda decyzja zapada po zebraniu kworum ponad dwóch trzecich głosów i jest natychmiast ostateczna.

Tendermint jest protokołem konsensusu bizantyjskiego wraz z implementacją silnika konsensusu. Jego pełną specyfikację oraz dowody bezpieczeństwa i żywotności zawiera praca \cite{buchman2018tendermint}. U jego podstaw leży spostrzeżenie, że klasyczne protokoły BFT rozwiązują ten sam problem, co kosztowne wydobycie w sieciach opartych na dowodzie pracy. Jest nim uzgodnienie wspólnej kolejności operacji wśród wzajemnie nieufających sobie uczestników. Zamiast konkurencji o prawo do zaproponowania bloku Tendermint wprowadza zatem jawny zbiór walidatorów głosujących nad kolejnymi blokami, przez co zatwierdzenie bloku jest decyzją kworum, a nie hipotezą podlegającą późniejszej rewizji.

Dla niniejszej pracy istotne są dwa aspekty tej konstrukcji. Pierwszym jest deterministyczna finalność, wynikająca wprost z głosowanego charakteru protokołu. Drugim jest rozdzielenie odpowiedzialności między silnik konsensusu a aplikację. Implementacja dostarcza kompletną warstwę konsensusu, obejmującą sieć P2P, pulę transakcji oczekujących (\emph{mempool}) oraz replikację bloków, natomiast to, co transakcje znaczą, pozostawia aplikacji komunikującej się z silnikiem przez wyspecyfikowany interfejs (sekcja~\ref{sec:abci}). Z tego powodu Tendermint bywa opisywany nie jako blockchain o konkretnym przeznaczeniu, lecz jako uniwersalny silnik replikacji dowolnej deterministycznej maszyny stanów. Protokół stanowi podstawę ekosystemu Cosmos oraz działających w nim sieci produkcyjnych, a jego utrzymywana obecnie implementacja nosi nazwę CometBFT \cite{cometbft-docs} i została wykorzystana w prototypie opisanym w dalszej części pracy.

Mechanika protokołu jest następująca. Znany zbiór walidatorów uzgadnia kolejne bloki w rundach złożonych z trzech kroków: wyznaczony rotacyjnie proponent rozgłasza blok (\emph{propose}), po czym walidatorzy głosują w dwóch turach (\emph{prevote}, \emph{precommit}); blok zostaje zatwierdzony, gdy ponad $2/3$ łącznej mocy głosów złoży \emph{precommit} w tej samej rundzie. Jeżeli kworum nie zostanie osiągnięte (awaria proponenta, opóźnienia sieci), protokół przechodzi do kolejnej rundy z nowym proponentem. Z konstrukcji wynikają trzy właściwości kluczowe dla niniejszej pracy:
\begin{itemize}
    \item \textbf{finalność natychmiastowa i deterministyczna} - zatwierdzony blok nie może zostać wycofany, dopóki złośliwi walidatorzy kontrolują mniej niż $1/3$ mocy głosów; rozwidlenia łańcucha nie występują, a ataki dalekiego zasięgu są bezprzedmiotowe;
    \item \textbf{bezpieczeństwo ponad żywotnością} (\emph{safety over liveness}) - gdy dostępne jest mniej niż $2/3$ mocy głosów, sieć wstrzymuje produkcję bloków zamiast ryzykować rozszczepienie historii; dla rejestru certyfikatów oznacza to, że awaria może czasowo uniemożliwić rejestrację nowych wpisów, ale nigdy nie doprowadzi do współistnienia dwóch sprzecznych widoków rejestru;
    \item \textbf{jawny zbiór walidatorów} - odpowiedzialność za każdy blok jest kryptograficznie przypisana do konkretnych, znanych podmiotów, co odpowiada konsorcyjnemu modelowi zaufania zarysowanemu w sekcji~\ref{sec:rozproszony-rejestr}.
\end{itemize}

Kosztem tych gwarancji jest komunikacja $O(n^2)$ pomiędzy walidatorami w każdej rundzie, ograniczająca praktyczny rozmiar zbioru walidatorów do rzędu setek. W zamierzonym zastosowaniu jest to ograniczenie akceptowalne, ponieważ liczba podmiotów tworzących konsorcjum rejestru i tak byłaby tego rzędu. Empiryczna demonstracja zachowania protokołu w poszczególnych scenariuszach awarii (tolerancja $f$ uszkodzeń, zatrzymanie poniżej kworum, odzyskiwanie) na uruchomionej sieci testowej stanowi część rozdziału~\ref{ch:analiza-bezpieczenstwa}.

\section{Interfejs ABCI i model aplikacji}
\label{sec:abci}

W niniejszej pracy wykorzystano CometBFT \cite{cometbft-docs}, utrzymywaną obecnie implementację protokołu Tendermint, stanowiącą kontynuację projektu Tendermint Core. Architektonicznym wyróżnikiem CometBFT, decydującym o jego wyborze, jest interfejs ABCI, który rozdziela system na dwie warstwy: silnik konsensusu, odpowiadający za sieć P2P, pulę transakcji oczekujących, uzgadnianie i replikację bloków, oraz aplikację, definiującą, co transakcje znaczą, czyli format danych, reguły walidacji i przejścia stanu. Aplikacja jest deterministyczną maszyną stanów komunikującą się z silnikiem przez wąski, jawnie wyspecyfikowany zestaw wywołań; cała logika DPKI opisana w kolejnych rozdziałach mieści się w tej warstwie, bez modyfikowania protokołu konsensusu.

Najważniejsze wywołania ABCI w wersji 2.0 dzielą się na trzy grupy. Na etapie przyjmowania transakcji \texttt{CheckTx} pełni rolę bramki puli transakcji oczekujących: pozwala tanio odrzucić transakcje nieprawidłowe, zanim trafią do propozycji bloku, przy czym jego werdykt nie jest częścią konsensusu. Na etapie uzgadniania \texttt{PrepareProposal} i \texttt{ProcessProposal} dają aplikacji wpływ odpowiednio na zawartość proponowanego bloku oraz na decyzję o głosowaniu nad cudzą propozycją. Wreszcie na etapie wykonania \texttt{FinalizeBlock} dostarcza aplikacji zatwierdzony blok do deterministycznego przetworzenia, a \texttt{Commit} utrwala wynikowy stan. Wywołania zapytaniowe (\texttt{Query}), obsługiwane poza ścieżką konsensusu, dopełniają interfejs o odczyt stanu.

Dwa elementy tego modelu mają bezpośrednie znaczenie dla konstrukcji systemu. Po pierwsze, aplikacja zwraca po wykonaniu każdego bloku skrót swojego stanu (\emph{app hash}), który silnik konsensusu umieszcza w nagłówku bloku \emph{następnego}. Tym samym stan aplikacji po wysokości $H$ jest kryptograficznie poświadczony podpisami walidatorów w nagłówku bloku $H+1$. Jeżeli skrótem stanu aplikacji jest korzeń rzadkiego drzewa Merkle (sekcja~\ref{sec:merkle}), każda odpowiedź na zapytanie o stan może zostać opatrzona dowodem weryfikowalnym względem zatwierdzonego nagłówka, bez zaufania do odpytywanego węzła. Po drugie, rozdzielenie \texttt{CheckTx} od \texttt{FinalizeBlock} wyznacza granicę determinizmu: walidacja w \texttt{FinalizeBlock} musi być w pełni deterministyczna (każda rozbieżność rozszczepiłaby skróty stanów aplikacji walidatorów), natomiast \texttt{CheckTx} i \texttt{ProcessProposal}, działające przed konsensusem, mogą sięgać do źródeł zewnętrznych. Rozróżnienie to okaże się kluczowe przy projektowaniu weryfikacji własności domeny (rozdział~\ref{ch:projekt}), wymagającej z natury niedeterministycznych zapytań DNS.
